\documentclass[problemstyle=boxed,theme=color,captionstyle=compact,subnumbering,tightlists,mathdelims,statmacros,localizedate]{homework}

% Recommended configuration: line spread, margins, heading spacing, solution spacing
\homeworksetlinespread{1.12}
\homeworksetmargins{left=1in,right=1in,top=0.95in,bottom=1.05in}
\homeworksetproblemtitlespacing{*1.6}{*1.0}
\homeworksetsolutionskip{0.55\baselineskip}{0.45\baselineskip}

% Metadata and localized labels
\renewcommand{\hmwkTitle}{Assignment \#1 Demo}
\renewcommand{\hmwkDueDate}{March 15, 2025 (Saturday)}
\renewcommand{\hmwkClass}{Probability and Statistics}
\renewcommand{\hmwkClassTime}{Wed 3-4}
\renewcommand{\hmwkClassInstructor}{Dr. Zhang}
\renewcommand{\hmwkAuthorName}{Alice Smith}
\renewcommand{\hmwkSolutionName}{Solution Outline}
\renewcommand{\qedsymbol}{\ensuremath{\blacksquare}}

% Hooks: add tips and separators around each problem
\renewcommand{\homeworkBeforeProblemHook}{
    \par\small\textbf{Heads-up:} Review the task before writing your solution.\par\medskip
}
\renewcommand{\homeworkAfterProblemHook}{
    \par\medskip\rule{\linewidth}{0.4pt}\par\medskip
}

\begin{document}

\maketitle

\pagebreak

\begin{problem}
    \textbf{Learning objectives overview.} This problem showcases the essentials of the `homework` class:
    \begin{itemize}
        \item Automatic headers/footers and boxed problems via `problemstyle=boxed`.
        \item Compact lists (`tightlists`) with custom spacing helpers.
        \item Paired delimiters and probability macros from `mathdelims` and `statmacros`.
    \end{itemize}

    Assume random variables $X$ and $Y$ satisfy $\Expect{X}=0$, $\Var{X}=1$, and $\Cov(X,Y)=\rho$. Complete the subtasks:
    \begin{enumerate}
        \item Prove that $\abs{\Cov(X,Y)} \leq \norm{Y}_2$.
        \item Derive bounds for $\Prob{\abs{Y}\leq 1}$ and simplify $\indicator{Y>0}$.
        \item Provide the formula for $\BiasOf{\hat\theta}$ and decide whether it vanishes.
    \end{enumerate}

    \solution

    \begin{enumerate}
        \item By Cauchy-Schwarz, $\abs{\Cov(X,Y)}=\abs{\Expect{XY}}\leq\norm{X}_2\norm{Y}_2=\norm{Y}_2$.
        \item Note $\Prob{\abs{Y}\leq 1}=1-\Prob{\abs{Y}>1}$, and Markov gives $\Prob{\abs{Y}>1}\leq\Expect{\abs{Y}}$. Also $\Expect{\indicator{Y>0}}=\Prob{Y>0}$.
        \item For an estimator $\hat\theta$, $\BiasOf{\hat\theta}=\Expect{\hat\theta}-\theta$. If $\Expect{\hat\theta}=\theta$, the bias is zero.
    \end{enumerate}

    \begin{note}[Further reading]
        Add references or textbook sections here when students need extra context.
    \end{note}
\end{problem}

\pagebreak

\begin{problem}[Method of Moments vs. MLE]
    Suppose samples $\{X_i\}_{i=1}^n$ and $\{Y_i\}_{i=1}^n$ come from the same parametric family with parameter $\theta$. Use techniques from class.

    \begin{subproblem}
        Write the moment estimator and comment on unbiasedness.
    \end{subproblem}
    \begin{subproblem}
        Derive the likelihood equation and discuss initial guesses for Newton iterations.
    \end{subproblem}
    \begin{subproblem*}[Discussion]
        Compare estimator variances and explain when one is more robust.
    \end{subproblem*}

    \solution

    \begin{enumerate}
        \item Equate sample moments to theoretical moments; linear models often yield unbiased estimators.
        \item Solve $\partial\ell(\theta)/\partial\theta=0$ and start Newton iterations at $\hat\theta_{\mathrm{MoM}}$.
        \item When data include outliers, prefer the estimator with lower variance or apply robust tweaks.
    \end{enumerate}
\end{problem}

\pagebreak

\homeworksetproblemtitle{Supplemental: Adaptive Layout}
\begin{problem*}
    This task demonstrates overriding the upcoming title and freezing numbering via `problem*`.
    Follow the steps below:
    \begin{enumerate}
        \item Show that if $f\colon \mathbb{R}\to\mathbb{R}$ is differentiable with $f(0)=0$, then there exists $\xi\in(0,1)$ such that $f(1)=f'(\xi)$.
        \item Use the `proof` environment so the $\blacksquare$ appears automatically at the end.
    \end{enumerate}

    \solution

    \begin{proof}
        Define $g(x)=f(x)-f(0)-f'(0)x$. Then $g(0)=g(1)=0$. By the mean value theorem, there exists $\xi\in(0,1)$ with $g'(\xi)=0$, which yields $f(1)=f'(\xi)$.
    \end{proof}

    \begin{note}
        To further customize layout, combine `homeworkBeforeProblemHook` and `homeworkAfterProblemHook` with additional structured material.
    \end{note}
\end{problem*}

\end{document}
